\section*{Chapter \ref{ch:nw:time-synchronisation} --- Time Synchronisation}

\begin{solution}{exo:nw:time-synchronisation:1}
It gains \qty{15}{\micro\second} a second: $100/15 = \qty{6.7}{\second}$ to breach 100 microseconds, $50\,000/15 = 3\,333$~s, about 56
minutes, to breach 50 milliseconds.
\end{solution}

\begin{solution}{exo:nw:time-synchronisation:2}
$\hat\vartheta = \tfrac12[(5\,300 - 0) - (8\,900 - 6\,300)] = \qty{1\,350}{\nano\second}$; round trip $8\,900 - 1\,000 =
\qty{7\,900}{\nano\second}$.
\end{solution}

\begin{solution}{exo:nw:time-synchronisation:3}
The asymmetry is $40 \times 4.88 = \qty{195.1}{\nano\second}$; half of it, \qty{97.5}{\nano\second}, remains in every estimate.
\end{solution}

\begin{solution}{exo:nw:time-synchronisation:4}
The timestamps' noise (\qty{5}{\micro\second} in software, nanoseconds in hardware) is small against queueing of several
microseconds per switch in each direction, which the estimate inherits as a random asymmetry. The least-delayed exchange of eight is
the one that queued least in both directions, so its asymmetry is smallest.
\end{solution}

\begin{solution}{exo:nw:time-synchronisation:5}
Solve $\tfrac12 \times 10^{-5} t^2 + t = 100\,000$: $t = (\sqrt{1 + 2} - 1)/10^{-5} = 73\,205$~s, about 20.3 hours.
\end{solution}

\begin{solution}{exo:nw:time-synchronisation:6}
Below a few seconds (the curves cross between 2 and 8~s): the servo's corrections jitter the disciplined clock by nanoseconds from one exchange to the next, while the free
oscillator's rate is steady over short times. What matters for the rule is the offset, which the servo keeps at nanoseconds; over
longer times the free clock's wandering rate dominates, and there the disciplined clock is better by orders of magnitude.
\end{solution}

\begin{solution}{exo:nw:time-synchronisation:7}
With \qty{50}{\micro\second} of queueing per switch and direction: about \qty{259}{\micro\second} with hardware timestamps alone and
\qty{111}{\micro\second} with the min-delay filter, both beyond the rule; \qty{25}{\nano\second} with \omterm{def:nw:time-synchronisation:bc}{transparent clocks}, which remove
queueing whatever its size.
\end{solution}

\begin{solution}{exo:nw:time-synchronisation:8}
PTP statistics measure the offset to the master along the path, with the path's fixed asymmetry invisible; they say nothing about the
master's own error to UTC, the card-to-host step, or where the timestamp is really taken. The 50~ns is one link of the chain, not the
chain.
\end{solution}

\begin{solution}{pb:nw:time-synchronisation:1}
\begin{enumerate}
\item $30 + 20 + 40 + 50 + 50 + 100 + 20 + 500 + 1\,500 = \qty{2\,310}{\nano\second}$.
\item $\sqrt{2\,518\,300} \approx \qty{1\,587}{\nano\second}$.
\item The software timestamp point and the card-to-host step: $2\,000/2\,310$, 87\%.
\item $100\,000 - 2\,310 = \qty{97\,690}{\nano\second}$.
\item Worst case \qty{310}{\nano\second}, root-sum-square $\sqrt{18\,300} \approx \qty{135}{\nano\second}$.
\item The card stamps the packet with its own clock, which PTP disciplines directly; the host's clock is no longer involved.
\item The residual \omterm{def:nw:time-synchronisation:asym}{path asymmetry}, \qty{100}{\nano\second}.
\item Calibrating the asymmetry away and removing a \omterm{def:nw:time-synchronisation:bc}{boundary clock}.
\item Solve $\tfrac12 \times 10^{-5} t^2 + t = 97\,690$: about 71\,866~s, 20.0 hours.
\item Solve $\tfrac12 \times 10^{-3} t^2 + 50\,t = 97\,690$: about 1\,917~s, 32 minutes.
\item Neither: 60 hours is three times the better one's \omterm{def:nw:time-synchronisation:holdover}{holdover}.
\item Record every minute in \omterm{def:nw:time-synchronisation:holdover}{holdover}, alarm at once, and switch to a second, independent reference if one exists.
\item \textbf{Named result.} The chain's worst case is \qty{2\,310}{\nano\second} with software timestamps and \qty{310}{\nano\second} with
  card timestamps; the grandmaster holds the remaining margin for 20.0 hours with the good oscillator and 32 minutes with the cheap one.
\item Timestamps must resolve a microsecond: a software timestamp read from a microsecond clock satisfies it, but its point of
  application must also be identified and consistent.
\item \omterm{def:nw:time-synchronisation:holdover}{Holdover} is bounded however good the oscillator is; an independent reference keeps the chain traceable, and lets each reference be
  checked against the other.
\item The documented design, the exact timestamp points, continuous monitoring records and an annual review.
\item The satellite's broadcast offset and a calibrated asymmetry, measured against an independent receiver or a portable reference.
\item Add \qty{10}{\micro\second}, more than the rest of the chain together, still inside the rule.
\item The American 50-millisecond limit is more than twenty thousand times the chain's worst case.
\item Every link from UTC to the point where the timestamp is taken, with its uncertainty.
\end{enumerate}
\end{solution}

\begin{solution}{iq:nw:time-synchronisation:1}
$t_2 = t_1 + d_{\mathrm{ms}} + \vartheta$, $t_4 = t_3 - \vartheta + d_{\mathrm{sm}}$, so
$\tfrac12[(t_2 - t_1) - (t_4 - t_3)] = \vartheta + \tfrac12(d_{\mathrm{ms}} - d_{\mathrm{sm}})$. It assumes equal delays both ways; the
asymmetry cannot be measured by exchanges.
\iqlookfor{the derivation and the unmeasurable asymmetry.}
\end{solution}

\begin{solution}{iq:nw:time-synchronisation:2}
A \omterm{def:nw:time-synchronisation:bc}{boundary clock} terminates PTP: it is a slave upstream and a master downstream. A \omterm{def:nw:time-synchronisation:bc}{transparent clock} forwards PTP and corrects each
message for its residence time. Either removes the switch's queueing from the estimate, which is otherwise the largest error on a
loaded trading network.
\iqlookfor{queueing as the problem, and the two ways of removing it.}
\end{solution}

\begin{solution}{iq:nw:time-synchronisation:3}
PTP with hardware timestamps and participating switches reaches nanoseconds to tens of nanoseconds on a local network; NTP with
software timestamps reaches microseconds at best, milliseconds over a wide area. NTP is enough for hosts whose timestamps must meet a
millisecond rule, such as a 50-millisecond audit-trail requirement.
\iqlookfor{hardware timestamps, participating switches, and matching the tool to the rule.}
\end{solution}

\begin{solution}{iq:nw:time-synchronisation:4}
The two clocks' synchronisation logs at that moment (offsets, \omterm{def:nw:time-synchronisation:holdover}{holdover}, alarms), where each timestamp was taken (card or software, send or
receive), any fixed asymmetry of the paths to the grandmaster, and whether the difference is within the sum of both budgets; then look
for a genuine bug in the ordering.
\iqlookfor{error budgets before conclusions.}
\end{solution}

\begin{solution}{iq:nw:time-synchronisation:5}
Show the documented chain and budget, the monitoring records for that day (grandmaster lock, every slave's offset and delay, no
\omterm{def:nw:time-synchronisation:holdover}{holdover} or excursions), the calibration of fixed asymmetries, the timestamp points, and the last annual review.
\iqlookfor{records and traceability, not a single number.}
\end{solution}

\begin{solution}{iq:nw:time-synchronisation:6}
The grandmaster goes into \omterm{def:nw:time-synchronisation:holdover}{holdover} and its offset grows with its oscillator's residual error and ageing; the slaves follow it. Switch to
an independent reference if there is one, alarm, compute from the oscillator's figures when the budget runs out, record the whole
period, and repair before then.
\iqlookfor{\omterm{def:nw:time-synchronisation:holdover}{holdover} arithmetic and a second reference.}
\end{solution}
